\documentclass[a4paper,twoside]{article}

\usepackage{morewrites}

\usepackage[svgnames,dvipsnames,x11names]{xcolor}
\usepackage{graphicx}
\usepackage[english]{babel}
\usepackage[colorlinks=true, linkcolor=NavyBlue, urlcolor=NavyBlue, citecolor=NavyBlue]{hyperref}
\usepackage[a4paper]{geometry}
\usepackage{ts-typesetting}
\usepackage{ts-fonts}
\usepackage{xcolor}
\usepackage{epigraph}
\usepackage{marginnote}
\usepackage{multicol}
\usepackage{progressbar}
\usepackage{graphicx}
\usepackage{tcolorbox}
\usepackage{fvextra}
\usepackage{amsmath}
\usepackage{float}
\usepackage{seqsplit}
\renewcommand*{\marginfont}{
\footnotesize
\sloppy
\emergencystretch=1em
\hyphenpenalty=50
\exhyphenpenalty=50
}
\newlength{\tsmarginparavail}
\newlength{\tsmarginparalt}
\newlength{\tsmarginparbuf}
\setlength{\tsmarginparbuf}{6mm}
\makeatletter
\AtBeginDocument{
\setlength{\tsmarginparavail}{\dimexpr\paperwidth-\textwidth-1in-\oddsidemargin-\marginparsep-\tsmarginparbuf\relax}
\ifdim\tsmarginparavail<\z@\setlength{\tsmarginparavail}{\z@}\fi
\if@twoside
\setlength{\tsmarginparalt}{\dimexpr1in+\evensidemargin-\marginparsep-\tsmarginparbuf\relax}
\ifdim\tsmarginparalt<\z@\setlength{\tsmarginparalt}{\z@}\fi
\ifdim\tsmarginparalt<\tsmarginparavail
\setlength{\tsmarginparavail}{\tsmarginparalt}
\fi
\fi
\ifdim\tsmarginparavail<\marginparwidth
\setlength{\marginparwidth}{\tsmarginparavail}
\fi
}
\makeatother
\usepackage{ts-code}

\usepackage{lastpage}
\usepackage{refcount}
\makeatletter
\def\ps@plain{
\let\@oddhead\@empty
\let\@evenhead\@empty
\def\@oddfoot{
\hfil
\ifnum\getpagerefnumber{LastPage}>1\relax
\thepage
\fi
\hfil}
\let\@evenfoot\@oddfoot
}
\makeatother

\title{TeX Engines}
\author{}\date{}
\begin{document}
\maketitle
\TeX{} has grown far beyond Knuth's original engine, evolving into a whole ecosystem of specialized typesetting machines. Each engine inherits the soul of classic \TeX{} but adds its own twist---some focusing on programmability, others on Unicode, scripting, or a modern toolchain experience. Together they form a strange but delightful family tree where 1980s design meets cutting-edge typography.
\section{\TeX{}}
The original \TeX{} engine, created by Donald Knuth, is the minimalist mathematical core of the entire ecosystem. It's deterministic, stable to the point of obsession, and designed so its output will match \emph{forever}. It handles typesetting with surgical precision but offers no frills---no Unicode, no PDF output, and no modern scripting hooks. Pure, legendary, and a little bit stubborn.
\section{E-\TeX{}}
e-\TeX{} extends \TeX{} with much-needed programming features without altering the underlying output. It adds new registers, improved conditionals, and tracing tools, making it a favorite for macro designers and format creators (like \LaTeX{}). Think of it as \TeX{} with a Swiss-army-knife upgrade.
\section{\tslogo{pdfTeX}}
\tslogo{pdfTeX} brought \TeX{} into the era of digital documents by producing PDF natively instead of going through DVI. It introduced microtypography---character protrusion, font expansion, and other subtle magic that makes text look professionally polished. Most modern \LaTeX{} distributions still rely heavily on \tslogo{pdfTeX}.
\section{\tslogo{XeTeX}}
\tslogo{XeTeX} is the engine that finally made \TeX{} feel Unicode-native. It uses system fonts directly (TrueType, OpenType), supports complex scripts naturally, and works beautifully for multilingual documents. If you need Arabic, Chinese, Hindi, or emoji without pain, \tslogo{XeTeX} is your friend.
\section{\tslogo{LuaTeX}}
\tslogo{LuaTeX} embeds a full Lua interpreter into the engine, effectively giving \TeX{} a programmable runtime. This allows deep customization, dynamic content generation, and powerful extensions like \tscodeinline{luaotfload} and \tscodeinline{luametalatex}. It's the most flexible and hackable \TeX{} engine, almost a \TeX{}/Lua hybrid organism.
\section{Tectonic}
Tectonic is a modern, Rust-powered \TeX{} engine aiming for reproducibility and user-friendliness. It automatically fetches missing packages, builds in a sandboxed environment, and removes the traditional ``\TeX{} installation anxiety.'' It tries to make \TeX{} behave like a modern build tool with zero configuration.
\section{Omega ($\Omega$) / Aleph ($\aleph$)}
Omega (and its successor Aleph) were early attempts at adding Unicode and advanced multilingual typesetting. They never became mainstream, but their ideas paved the way for \tslogo{XeTeX} and \tslogo{LuaTeX}.
\section{pTeX / upTeX}
Specialized engines designed for Japanese typesetting. pTeX handles vertical writing and Japanese line-breaking rules, while upTeX brings Unicode support to that world. They're essential in the Japanese \TeX{} community.
\section{Which to prefer?}
That a debate as old as \TeX{} itself still rages on is a testament to its complexity and versatility. For most users, \textbf{\tslogo{pdfTeX}} or \textbf{\tslogo{LuaTeX}} (with \LaTeX{} macros) will cover nearly all needs. Some facts:
\begin{enumerate}
\item Tectonic is so smooth, it downloads packages automatically, making it great for newcomers. No need to install a heavy \TeX{} distribution.
\item \tslogo{LuaLaTeX} is the only engine that supports both protrusion and font expansion (microtypography) along with Lua scripting, which yields smoother PDF output.
\item \tslogo{XeLaTeX} has similar results to Tectonic but allows \tscodeinline{-\allowbreak{}-\allowbreak{}shell-\allowbreak{}escape} for minted code highlighting.
\end{enumerate}
\end{document}
